\begin{drillsolution}{d:24}\textbf{从前顶点计算第一次折射}\par\nobreak
图\ref{fig:d24q}中光从左向右。采用实物距为正、沿出射方向的像距为正、球心在出射侧的曲率半径为正的约定，第一面 $R_{1s}=+10$，长度以厘米计：
\[
\frac1{15}+\frac{1.50}{v_1}=\frac{1.50-1}{10},
\qquad v_1=-90.0.
\]
所以中间像在 $O_1$ 左侧 $90.0\,\mathrm{cm}$，第一面放大率
$\beta_1=-n_1v_1/(n_2u_1)=+4$。它是第一面的正立虚像，但对后表面来说，入射光实际从这个位置发散。

\par\medskip\Needspace{4\baselineskip}\noindent\textbf{用两顶点间距更新第二面的物距}\par\nobreak
图\ref{fig:d24s}以 $O_1$ 为位置零点，$O_2$ 的位置为 $20$，中间像位置为 $-90$，因此第二面的实物距为
\[
u_2=20-(-90)=110.0\,\mathrm{cm}.
\]
第二面球心在入射侧，$R_{2s}=-10$；其入射介质为玻璃、出射介质为空气，故
\[
\frac{1.50}{110}+\frac1{v_2}=\frac{1-1.50}{-10},
\qquad \boxed{v_2=27.5\,\mathrm{cm}}.
\]
\fig{fig:d24s}{按有向位置连接两次成像，左侧长距离作压缩示意}{\begin{tikzpicture}
\draw[axis](-.4,0)--(8,0)node[right]{$x/\mathrm{cm}$};
\foreach \x/\lab/\val in{0/A_1/-90,3/O_1/0,4.5/O_2/20,6.6/A'/47.5}{\fill(\x,0)circle(1.4pt);\node[above]at(\x,0){$\lab$};\node[below]at(\x,0){$\val$};}
\draw[dim](0,1)--(4.5,1)node[midway,above]{$u_2=110$};
\draw[dim](4.5,-1)--(6.6,-1)node[midway,below]{$v_2=27.5$};
\draw[world](1.4,-.1)--(1.55,.1)(1.65,-.1)--(1.8,.1);
\end{tikzpicture}}

\par\medskip\Needspace{4\baselineskip}\noindent\textbf{求总放大率并核对像的位置}\par\nobreak
第二面 $\beta_2=-1.50\times27.5/110=-0.375$，总放大率为
\[
\boxed{\beta=\beta_1\beta_2=-1.50}.
\]
最终像在 $O_2$ 右侧 $27.5\,\mathrm{cm}$，即图\ref{fig:d24s}中 $x=47.5\,\mathrm{cm}$，为倒立放大实像。两面的顶点相距 $20\,\mathrm{cm}$，这段传播决定了 $u_2$，不能把实心球直接当成厚度为零的薄透镜。
\end{drillsolution}

\begin{drillsolution}{d:25}\textbf{用球面几何保留有限入射角}\par\nobreak
按图\ref{fig:d25q}取 $O=(0,0)$、$C=(-R,0)$。令半径 $CP$ 与正轴夹角为 $\theta$，则
\[
h=R\sin\theta,\qquad P=(R\cos\theta-R,R\sin\theta).
\]
入射方向为 $(1,0)$，单位法线为 $(\cos\theta,\sin\theta)$。由反射定律，出射方向为
\[
(1,0)-2\cos\theta(\cos\theta,\sin\theta)
=(-\cos2\theta,-\sin2\theta).
\]
题设 $\theta<45^\circ$，所以反射光向左下方传播，与主轴相交。

\par\medskip\Needspace{4\baselineskip}\noindent\textbf{求实际反射线与主轴的交点}\par\nobreak
反射线从 $P$ 到主轴的路程为 $h/\sin2\theta=R/(2\cos\theta)$。其交点横坐标为
\[
x_F=R(\cos\theta-1)-\frac{R\cos2\theta}{2\cos\theta}
=\frac R{2\cos\theta}-R.
\]
于是图\ref{fig:d25s}中 $s=-x_F$，得到
\[
\boxed{s(h)=R-\frac{R}{2\sqrt{1-h^2/R^2}}}.
\]
\fig{fig:d25s}{有限高度光线的交点比近轴焦点更靠近顶点}{\begin{tikzpicture}
\draw[axisline](-4.3,0)--(.6,0);\draw[world](-3,0)++(-40:3)arc(-40:40:3);
\coordinate(P)at(-.401924,1.5);\coordinate(F)at(-1.267949,0);
\draw[ray](-4,1.5)--(P);\draw[ray](P)--(F);\draw[normal](-3,0)--(P);
\fill(-3,0)circle(1.5pt)node[below]{$C$};\node[right]at(0,0){$O$};\fill(F)circle(1.5pt);
\node at(-.85,-.45){$F(h)$};\draw[aux](F)--(-.85,-.25);
\fill(-1.5,0)circle(1.2pt)node[below left]{$F_0$};
\draw[aux](F)--(-1.267949,-2.3);\draw[aux](0,0)--(0,-2.3);
\draw[dim](-1.267949,-2.3)--(0,-2.3)node[midway,below]{$s(h)$};
\draw(-2.35,0)arc(0:30:.65);\node at(-2.08,.25){$\theta$};
\node at(-2.6,-1){$OF_0=R/2$};
\end{tikzpicture}}

\par\medskip\Needspace{4\baselineskip}\noindent\textbf{由焦点偏移确定孔径限制}\par\nobreak
当 $h/R=1/2$，$\cos\theta=\sqrt3/2$，故 $s=20(1-1/\sqrt3)\,\mathrm{cm}\simeq8.453\,\mathrm{cm}$；近轴值为 $10.0\,\mathrm{cm}$，相对偏差约 $15.47\%$。图\ref{fig:d25s}的交点顺序与 $s<R/2$ 一致。

一般地，相对偏差为 $\sec\theta-1$。要求它不超过 $0.0100$，等价于 $\cos\theta\ge100/101$，从而
\[
\boxed{h_{\max}=R\frac{\sqrt{201}}{101}\simeq2.807\,\mathrm{cm}}
\quad(R=20.0\,\mathrm{cm}).
\]
$h\to0$ 时精确式恢复 $s\to R/2$，而误差首项为 $h^2/(2R^2)$。孔径限制来自球面反射的几何像差，与衍射限制是不同的问题。
\end{drillsolution}
